% =====================================================================
%  13. Modern computational perspective
% =====================================================================
\section{Gradient methods today}

\begin{frame}{Why first-order methods won: gradients are cheap}
  \begin{thmbox}[Cheap gradient principle \normalfont\cite{griewank2008}]
    Reverse-mode automatic differentiation computes the \emph{whole} gradient $\nabla f\in\R^n$ for a small constant
    multiple of the cost of evaluating $f$ once, \alert{no matter how large $n$ is}.
  \end{thmbox}
  \vspace{2pt}
  \centering\small
  \begin{tabular}{@{}l c c@{}}
    \toprule
    How you get $\nabla f$ & Cost & Accuracy \\
    \midrule
    finite differences (\S9) & $(n+1)\times$ cost$(f)$ & $\sim\sqrt{\epsilon_{\text{mach}}}$ \\
    forward-mode AD & $\sim n\times$ cost$(f)$ & machine precision \\
    reverse-mode AD (backpropagation) & $O(1)\times$ cost$(f)$ & machine precision \\
    \bottomrule
  \end{tabular}

  \raggedright
  \begin{itemize}\footnotesize
    \item Backpropagation in neural networks \emph{is} reverse-mode AD. That is why the gradient of a $10^9$-parameter
          model costs about as much as a few forward passes.
    \item GD itself only needs matrix--vector style operations, which run in parallel on GPUs.
  \end{itemize}
\end{frame}

\begin{frame}{Adam, and a puzzle called the edge of stability}
  \begin{columns}[T]
    \begin{column}{0.5\textwidth}
      \textbf{Adam}~\cite{kingma2015}, with all operations element-wise:
      \small
      \begin{align*}
        \mathbf m_k &= \beta_1\mathbf m_{k-1}+(1-\beta_1)\g_k \\
        \vv_k &= \beta_2\vv_{k-1}+(1-\beta_2)\g_k^{2}\\
        \x_{k+1} &= \x_k-\alpha\,\hat{\mathbf m}_k/\bigl(\sqrt{\hat\vv_k}+\epsilon\bigr)
      \end{align*}
      \small Read it as \alert{momentum} plus a \alert{diagonal preconditioner} $\bP_k = \diag(\sqrt{\hat\vv_k})$. In the language
      of \S2, it is steepest descent in a diagonal norm that keeps changing. (The hats are a bias correction.)
    \end{column}
    \begin{column}{0.48\textwidth}
      \textbf{The edge of stability}~\cite{cohen2021}
      \begin{itemize}\small
        \item Classical theory says to keep $\alpha<2/L$.
        \item Yet in full-batch GD on neural networks, the \emph{sharpness} $\lmax(\nabla^2 f)$ is observed to climb until it
              reaches $2/\alpha$ and then hover just above it. The loss keeps going down over long horizons,
              though not monotonically.
        \item This is an \emph{empirical} finding that the classical analysis of \S5 to \S7 does not explain.
              It is an open research question.
      \end{itemize}
    \end{column}
  \end{columns}
\end{frame}
